%% kapa: exam=1 round=22 year=2011 period=1 form=A subject=civil_law subject_ko=민법 seq=36 number=36
%% source: Q-Net 감정평가사 기출문제 (cst00309) artlSeq=5036790 file=extracted/1008/5036790/2011년 제22회 감정평가사 1차시험 기출문제/1교시.hwp
%% conversion: kordoc 4.24.51 -> tools/gen1.py (auto); review=proofread-sonnet; flags=none
\begin{question}{36}
甲과 乙은 X토지를 공유하고 있는데, 甲의 지분은 2/3, 乙의 지분은 1/3이다. 다음 설명으로 옳은 것은? (다툼이 있으면 판례에 의함)
\begin{choices}
  \item 甲과 乙이 공유한 때로부터 10년 이상 경과하였다면 분할청구권이 소멸시효에 걸리므로, 乙은 甲에게 X토지의 분할을 청구할 수 없다.
  \item 乙이 관리비용 기타 의무이행을 1년 이상 지체한 때에는 甲은 상당한 가액으로 乙의 지분을 매수할 수 있다.
  \item 丙이 X토지를 불법으로 점유하고 있는 경우, 乙은 丙에 대해 단독으로 X토지 전체에 대해 손해배상을 청구할 수 있다.
  \item 甲ㆍ乙이 丙에게 X토지를 임대해 준 경우, 甲이 단독으로 계약을 해지하여도 계약은 그대로 존속한다.
  \item 甲이 乙의 동의 없이 X토지 전부를 丙에게 임대하여 丙이 점유하는 경우, 乙은 丙에게 단독으로 점유의 배제를 청구할 수 있다.
\end{choices}
\end{question}
\answer{2}
